\documentclass[12pt]{article}
\usepackage[utf8]{inputenc}
\usepackage[T1]{fontenc}
\usepackage{hyperref}
\title{Tavanlav Language Resources for K-12 Curriculum Development}
\author{Compiled by Claude Sonnet 5 (claude-sonnet-5) for Andr\'as Kornai}
\date{}
\begin{document}
\maketitle

\section{Introduction}
Tavanlav (ISO 639-3 vlp, Glottocode valp1237) is an Austronesian (Oceanic) language of Vanuatu. In my searches the name ``Tavanlav'' appeared only as a dialect or alternate name of the language usually called Valpei (also Valpei-Hukua). The sources I found place Valpei at the northern tip of Espiritu Santo. One Joshua Project page places the Valpei people on Pentecost, which conflicts with that. I did not resolve the conflict. Wikipedia gives about 300 speakers, citing 1983, while Joshua Project lists a population of 800. I found no school materials, primers or published dictionaries in these searches. Little appears to exist.

\section{Primary Educational Resources}
I found no primers, textbooks or school programmes for Valpei/Tavanlav.

\section{Linguistic and Reference Materials}
\begin{itemize}
\item Valpei language, Wikipedia overview: \url{https://en.wikipedia.org/wiki/Valpei_language}
\item Espiritu Santo languages, Wikipedia (lists Valpei, also called Tavanlav): \url{https://en.wikipedia.org/wiki/Espiritu_Santo_languages}
\item Glottolog language page (map view) for valp1237: \url{https://glottolog.org/resource/languoid/id/valp1237.bigmap.html}
\end{itemize}

\section{Religious and Community Resources}
\begin{itemize}
\item Joshua Project, language page for vlp (reports no confirmed complete Bible translation): \url{https://joshuaproject.net/languages/vlp}
\item Joshua Project, Valpei people group in Vanuatu: \url{https://joshuaproject.net/people_groups/15767}
\end{itemize}

\section{Digital and Audio Resources}
I found no audio, video or corpus resources for this language in these searches.

\section{Curriculum Development Context}
Documentation is very thin. Any curriculum would depend on new fieldwork and community-led work. The location discrepancy and the Valpei/Tavanlav naming should be settled with the community before building materials. The sources I found do not say whether the language is taught in schools.

\end{document}